% Neither End of the Bridge: terse theory paper for the Schur bridge.
% Self-contained. Build: tectonic neither-end-of-the-bridge.tex
% Certificate script: paper/neither_end_certificate.py (exact rational
% arithmetic; stdlib only).
\documentclass[11pt]{article}

\usepackage[utf8]{inputenc}
\usepackage[T1]{fontenc}
\usepackage{lmodern}
\usepackage[margin=1in]{geometry}
\usepackage{amsmath,amssymb,amsthm}
\usepackage{booktabs}
\usepackage{microtype}
\usepackage[hidelinks]{hyperref}

\newtheorem{theorem}{Theorem}
\newtheorem{lemma}{Lemma}
\newtheorem{proposition}{Proposition}
\theoremstyle{remark}
\newtheorem{remark}{Remark}

\newcommand{\E}{\mathbb{E}}
\DeclareMathOperator*{\argminop}{arg\,min}

\title{Neither End of the Bridge:\\ Estimation Noise Makes Interior Schur Coupling Optimal}
\author{Peter Cotton\thanks{[Affiliation and postal address to be supplied.]
Email: \texttt{peter.cotton@microprediction.com}. The machine-verifiable
certificate for Proposition~\ref{prop:witness} is
\texttt{paper/neither\_end\_certificate.py} in
\texttt{github.com/microprediction/schur}; implementations of the recursion are in
the \texttt{allocation} and \texttt{precise} packages.}}
\date{\today}

\begin{document}
\maketitle

\begin{abstract}
Schur-complementary allocation interpolates hierarchical risk parity (HRP,
$\gamma = 0$) and minimum variance ($\gamma \to 1$) through a coupling
$\gamma$ that scales how much of the estimated cross-block covariance the
recursion uses. Antonov, Lipton, and L\'opez de Prado prove that of the two
endpoints, HRP is the less noisy; Mograby shows the $\gamma = 1$ endpoint is
computable exactly on hierarchical structures. We prove that, generically,
neither endpoint is optimal. The $\gamma = 0$ recursion never reads the
top-level cross-block, so its expected out-of-sample variance $F(0)$ is exact
while $F'(0)$ equals the strictly negative population slope: HRP is never the
minimizer under small noise. Every appearance of the estimated cross-block
carries a factor of $\gamma$, so the noise cost is $O(\gamma^2 \tau^2)$ near
zero yet order one at $\gamma = 1$: above an explicit noise threshold, full
coupling is never the minimizer either. A four-asset witness makes the
conclusion unconditional --- $F(9/10) < F(0)$ and $F(9/10) < F(1)$ hold as
exact rational inequalities, machine-verifiable because the recursion is
rational arithmetic. Finally, under decreasing-convex population value and
increasing-convex noise cost, the optimal coupling is the unique root of
$\tau^2 = -V_0'(\gamma)/G'(\gamma)$ and sweeps the bridge monotonically:
every interior $\gamma$ is optimal at exactly one noise level, full coupling
only as noise vanishes, HRP only as it diverges.
\end{abstract}

\noindent\textbf{Keywords:} portfolio construction; hierarchical risk parity;
Schur complement; estimation error; shrinkage; bias--variance trade-off.\\
\textbf{JEL:} G11, C13, C58. \quad \textbf{MSC2020:} 91G10, 62H12.

\section{Introduction}
\label{sec:intro}

Hierarchical risk parity \cite{lopezdeprado2016} allocates by recursive
bisection without inverting a covariance matrix; minimum variance inverts it
in full. That the two are ends of a single construction --- augmenting each
block with a $\gamma$-scaled Schur complement of its sibling yields a
recursion that is exactly HRP at $\gamma = 0$ and recovers minimum variance
as $\gamma \to 1$ --- is a unification that first appeared in a blog
\cite{cotton2022blog}, then in a textbook \cite[\S12.3.4]{palomar2025}, with
a fuller development in \cite{cotton2024schur}; on the hierarchical graph
structures of \cite{mograby2025} the $\gamma = 1$ recovery is exact. The
endpoints have since been compared analytically: \cite{antonov2025} derive
the weight noise of HRP and of Markowitz under covariance estimation error
and prove HRP is the more robust.

This paper proves the statistical claim implicit in the bridge itself: under
estimation noise the optimal coupling is generically \emph{interior}, and its
location is characterized by a bias--variance first-order condition. The
contributions are four short results. (i)~A structural lemma: the
$\gamma = 0$ recursion never reads the top-level cross-block, so cross-block
noise cannot move it. (ii)~If coupling has first-order population value, HRP
is strictly suboptimal under all sufficiently small noise, of any support.
(iii)~Because the estimated cross-block enters the recursion only with a
$\gamma$ prefactor, the noise cost vanishes quadratically at the HRP end
while it is order one at full coupling; above an explicit threshold the
minimizer is interior --- and a four-asset witness, verified in exact
rational arithmetic, makes the conclusion unconditional. (iv)~Under shape
conditions verified in the witness family, the optimal coupling
$\gamma^*(\tau^2)$ is the unique root of $\tau^2 = -V_0'/G'$ and decreases
monotonically in the noise, sweeping the entire bridge: an optimal-shrinkage
formula in the spirit of \cite{ledoitwolf2004}, with HRP the infinite-noise
limit and full coupling the zero-noise limit.

\section{The recursion}
\label{sec:setup}

Fix an asset order and a balanced bisection tree over $n$ assets. Given
$M \succ 0$ and $\gamma \in [0,1]$, the Schur-complementary allocation
$w(\gamma; M)$ is defined recursively \cite{cotton2024schur}: at a node with
children $(L, R)$ and blocks
$M = \left(\begin{smallmatrix} A & B \\ B^\top & D \end{smallmatrix}\right)$,
each child block is augmented with the $\gamma$-scaled Schur complement of
the other,
\begin{equation}
  A_\gamma \;=\; \operatorname{sym}\!\Big[\big(I - \gamma B D^{-1}\big)^{-1}
  \big(A - \gamma\, B D^{-1} B^\top\big)\Big],
  \qquad
  D_\gamma \;=\; \big(A \leftrightarrow D,\ B \to B^\top\big),
  \label{eq:aug}
\end{equation}
the budget is split between children in inverse proportion to the naive
portfolio variances $\nu(A_\gamma), \nu(D_\gamma)$ (inverse-diagonal
weights), and the recursion continues inside each augmented block. At
$\gamma = 0$ the augmentation is the identity and the recursion is HRP over
the fixed order.

Let $\Sigma \succ 0$ be the population covariance and
$\widehat\Sigma_\tau = \Sigma + \tau N$ an estimate, with $N$ a bounded
random symmetric perturbation on a set where every matrix in \eqref{eq:aug}
remains invertible for $\gamma \in [0, \bar\gamma]$. The object of study is
the expected out-of-sample variance and its population limit,
\begin{equation}
  F(\gamma; \tau) \;=\; \E\big[\, w(\gamma; \widehat\Sigma_\tau)^\top \Sigma\,
  w(\gamma; \widehat\Sigma_\tau) \,\big],
  \qquad
  V_0(\gamma) \;=\; F(\gamma; 0).
  \label{eq:F}
\end{equation}
Every operation in the recursion is rational in $(\gamma, M)$, so $w$, $F$,
$V_0$ are rational functions where denominators do not vanish --- in
particular $C^\infty$ there, which is all the regularity used below.

\section{Neither endpoint}
\label{sec:endpoints}

\begin{lemma}[The HRP end is blind to the cross-block]
\label{lem:blind}
$w(0; M)$ is a function of $(A, D)$ alone. Consequently, if the estimation
error is supported on the top-level cross-block, then
$w(0; \widehat\Sigma_\tau) = w(0; \Sigma)$ almost surely and
$F(0; \tau) = V_0(0)$ exactly, for every $\tau$.
\end{lemma}
\begin{proof}
At $\gamma = 0$ the augmentation \eqref{eq:aug} returns its first argument,
so the top-level split uses $\nu(A)$ and $\nu(D)$, which read only
within-block entries, and the recursion continues inside the unmodified
blocks $A$ and $D$, never touching $B$. (Cross-blocks interior to $A$ or $D$
\emph{are} read, through the naive variances of their parents; only the
top-level $B$ is never read.)
\end{proof}

\begin{remark}[Scope]
\label{rem:scope}
The lemma is per-tree and per-variant: this recursion (naive variance over
full child sub-blocks, as in \cite{lopezdeprado2016}) never reads the
top-level cross-block; a split on diagonal variances alone would read nothing
off-diagonal. When the tree itself is data-driven (a dendrogram, or Fiedler
seriation \cite{cotton2024schur}), the ordering step does read the
cross-block --- but only through the discrete tree choice, which for a
generic estimate is locally constant. The $\gamma = 0$ allocation is then
locally constant in the cross-block entries, and the differential statements
below are unaffected for noise small enough to leave the tree unchanged. (In
the witness of Proposition~\ref{prop:witness} the seriation returns the same
top partition on both noise branches, so fixing the order there is without
loss.)
\end{remark}

\begin{theorem}[HRP is strictly suboptimal under small noise]
\label{thm:nohrp}
If $V_0'(0) < 0$ --- coupling has first-order population value --- then there
is $\tau_0 > 0$ such that for all $\tau < \tau_0$,
$\partial_\gamma F(0; \tau) < 0$: every minimizer of $F(\cdot\,;\tau)$ on
$[0, \bar\gamma]$ is bounded away from $0$. No support restriction on $N$ is
needed.
\end{theorem}
\begin{proof}
$\partial_\gamma F(0;\tau) = \E[\partial_\gamma V(0; \Sigma + \tau N)]$,
where $V(\gamma; M) = w(\gamma; M)^\top \Sigma\, w(\gamma; M)$ is $C^1$ near
$(0, \Sigma)$ with $\partial_\gamma V$ continuous and, for bounded $N$,
dominated there. By dominated convergence
$\partial_\gamma F(0;\tau) \to V_0'(0) < 0$ as $\tau \to 0$; continuity in
$\tau$ gives $\tau_0$.
\end{proof}

\begin{theorem}[Full coupling is suboptimal above a noise threshold]
\label{thm:nofull}
Let the error be supported on the top-level cross-block and symmetric
($N \overset{d}{=} -N$). Then:
\begin{enumerate}
\item[\textup{(i)}] there is $C < \infty$, uniform on a compact set of
$\gamma$ and noise realizations keeping \eqref{eq:aug} invertible, with
$\big| F(\gamma;\tau) - V_0(\gamma) \big| \le C \gamma^2 \tau^2$;
\item[\textup{(ii)}] if for some $\hat\gamma \in (0,1)$ and $\tau^*$ the
excess noise cost at full coupling satisfies
$F(1;\tau^*) - V_0(1) > V_0(\hat\gamma) - V_0(1) + C\hat\gamma^2\tau^{*2}$,
then $F(\hat\gamma;\tau^*) < F(1;\tau^*)$, and with
Theorem~\ref{thm:nohrp} every minimizer of $F(\cdot\,;\tau^*)$ is interior.
\end{enumerate}
\end{theorem}
\begin{proof}
(i) In \eqref{eq:aug} the cross-block enters only through
$\gamma \widehat B D^{-1}$ and $\gamma \widehat B D^{-1} \widehat B^\top$,
so every occurrence of the perturbation $\tau N_B$ carries the prefactor
$\gamma$, and the deeper levels operate on augmented blocks that retain it.
Hence $\partial w/\partial\tau = \gamma\, h(\gamma,\tau)$ with $h$ bounded on
the stated compact set, so $V(\gamma;\tau N) - V(\gamma;0)$ is Lipschitz in
$\tau$ with constant $O(\gamma)$; by symmetry of $N$ the first-order term in
$\tau$ cancels in expectation, and Taylor's theorem with the uniformly
bounded second derivative gives (i). (ii) is arithmetic:
$F(\hat\gamma;\tau^*) \le V_0(\hat\gamma) + C\hat\gamma^2\tau^{*2}
< F(1;\tau^*)$.
\end{proof}

\begin{remark}
The sign of $F - V_0$ is not fixed: symmetric noise can reduce expected
out-of-sample variance slightly at moderate $\gamma$ through the curvature of
$V$ in the perturbation (observed in the witness family). The theorems bound
magnitudes; what they use is only that noise costs nothing at $\gamma = 0$
and that its magnitude at $\gamma = 1$ can exceed the population gain of full
coupling --- which it does dramatically once the perturbed recursion
approaches the boundary of the positive-definite cone, where the
$\gamma = 1$ allocation is not defined at all.
\end{remark}

\begin{proposition}[Exact witness]
\label{prop:witness}
Let $n = 4$ with blocks $\{1,2\}, \{3,4\}$, volatilities $(1,1,2,2)$,
within-block correlation $\tfrac12$, cross-block correlation $\tfrac{3}{10}$;
let the estimate err in the $(1,3)$ covariance entry by
$\pm \tfrac{3}{20}\cdot 2$ with probability $\tfrac12$ each. Then, in exact
rational arithmetic,
\[
  F(\tfrac{9}{10}) < F(0) \ \ \text{(margin } 4.879\times10^{-2}\text{)},
  \qquad
  F(\tfrac{9}{10}) < F(1) \ \ \text{(margin } 5.702\times10^{-4}\text{)},
\]
so every minimizer of $F$ on $[0,1]$ is strictly interior. All augmented
blocks remain positive definite on both noise branches for every
$\gamma \in [0,1]$: no repair or projection is involved.
\end{proposition}
\begin{proof}
$F$ at rational arguments is an exact rational number, being a finite
composition of field operations \eqref{eq:aug} on rational inputs. The two
inequalities are evaluated in exact fraction arithmetic by the accompanying
certificate script (see the title footnote), which also certifies positive
definiteness along the path.
\end{proof}

\section{The bridge as a bias--variance frontier}
\label{sec:frontier}

Smoothness gives the uniform second-order expansion
\begin{equation}
  F(\gamma;\tau) \;=\; V_0(\gamma) + \tau^2 G(\gamma) + O(\tau^4),
  \qquad G(0) = G'(0) = 0,
  \label{eq:expansion}
\end{equation}
the boundary behavior of the noise-cost profile $G$ coming from
Lemma~\ref{lem:blind} and the $\gamma$-prefactor of
Theorem~\ref{thm:nofull}(i). Write $F_2 = V_0 + \tau^2 G$.

\begin{theorem}[The optimal coupling sweeps the bridge monotonically]
\label{thm:frontier}
Suppose on $(0,\bar\gamma)$ that $V_0' < 0$, $V_0'' \ge 0$ (decreasing convex
population curve) and $G' > 0$, $G'' \ge 0$ (increasing convex noise cost),
with at least one of $V_0'', G''$ strictly positive. Then
$t(\gamma) = -V_0'(\gamma)/G'(\gamma)$ is a continuous, strictly decreasing
bijection onto its range, and for every $\tau^2$ in that range the surrogate
$F_2(\cdot\,;\tau)$ has the unique minimizer
$\gamma^*(\tau^2) = t^{-1}(\tau^2)$, strictly decreasing in the noise level.
Within the surrogate: $\gamma^* \to \bar\gamma$ only as the noise vanishes,
$\gamma^* \to 0$ only as it diverges, and every interior coupling is optimal
at exactly one noise level.
\end{theorem}
\begin{proof}
$t' = (-V_0'' G' + V_0' G'')/G'^2 < 0$ on the stated domain, so $t$ is a
strictly decreasing bijection; $F_2' = V_0' + \tau^2 G'$ is negative before
$t^{-1}(\tau^2)$ and positive after. The comparative statics
$d\gamma^*/d\tau^2 = -G'/F_2'' < 0$ follow from the implicit function
theorem.
\end{proof}

\begin{proposition}[Transfer to the exact objective]
\label{prop:transfer}
Suppose in addition $F_2'' \ge c > 0$ on a neighborhood of
$\gamma^*(\tau^2)$ and $|F - F_2| \le K\tau^4$ there. Then any minimizer of
$F$ in that neighborhood satisfies
$|\gamma_F - \gamma^*| \le 2\sqrt{K/c}\;\tau^2$.
\end{proposition}
\begin{proof}
$\tfrac{c}{2}(\gamma_F - \gamma^*)^2 \le F_2(\gamma_F) - F_2(\gamma^*)
\le \big[F(\gamma_F) - F(\gamma^*)\big] + 2K\tau^4 \le 2K\tau^4$.
\end{proof}

\begin{remark}[Interpretation]
$\gamma$ is a shrinkage intensity toward the noise-immune end of the bridge,
and Theorem~\ref{thm:frontier} is an optimal-shrinkage formula in the spirit
of \cite{ledoitwolf2004}: the first-order condition
$\tau^2 = -V_0'(\gamma^*)/G'(\gamma^*)$ balances the marginal population
value of coupling against its marginal noise cost. The endpoint behaviors
recover the prior results. As $\gamma \to 1$ the augmentation approaches the
exact Schur complement and $G'$ inherits the inversion sensitivity of minimum
variance --- the instability quantified at the endpoint by
\cite{antonov2025} --- so full coupling is optimal only in the
vanishing-noise limit, exactness \cite{mograby2025} notwithstanding. As
$\gamma \to 0$, $G'$ vanishes while the population gain $-V_0'(0)$ does not,
so HRP is optimal only in the infinite-noise limit.
\end{remark}

\begin{remark}[Verification in the witness family]
\label{rem:verify}
With symmetric two-point noise on the cross-block correlation of the
Proposition~\ref{prop:witness} family, the measured profiles satisfy the
hypotheses: $V_0$ is decreasing convex on $[0,1)$; $G$ is increasing convex
with the predicted quadratic vanishing at $0$
($G \approx 0.003, 0.032, 0.138, 0.448$ at $\gamma = 0.2, 0.4, 0.6, 0.8$),
and is $\tau$-stable for $\tau \le 0.2$, confirming the $O(\tau^4)$ error.
The root of $V_0' + \tau^2 G' = 0$ tracks the exact minimizer of $F$:
\begin{center}
\begin{tabular}{lcccccc}
\toprule
$\tau$ & $0.10$ & $0.15$ & $0.20$ & $0.25$ & $0.30$ & $0.35$ \\
\midrule
exact $\argminop F$        & $0.85$ & $0.70$ & $0.65$ & $0.55$ & $0.50$ & $0.45$ \\
root of $V_0' + \tau^2 G'$ & $0.83$ & $0.74$ & $0.66$ & $0.59$ & $0.53$ & $0.49$ \\
\bottomrule
\end{tabular}
\end{center}
(grid resolution $0.05$): the optimal coupling falls monotonically from near
full coupling toward the middle of the bridge as the noise grows, as
Theorem~\ref{thm:frontier} requires.
\end{remark}

\section{Closing}

Exactness of full coupling does not confer statistical optimality, and the
noise-robustness of HRP does not confer it either: under covariance
estimation error the Schur coupling is a bias--variance dial whose optimal
setting is interior, decreasing in the noise, and characterized by
$\tau^2 = -V_0'/G'$. The empirical counterpart --- a walk-forward bridge on
Dow constituents with one-year training windows, minimum near
$\gamma \approx 0.55$ and robust to Ledoit--Wolf shrinkage of the inputs ---
is at \texttt{allocation.microprediction.org}, alongside the interactive
demonstrations.

\begin{thebibliography}{9}

\bibitem{lopezdeprado2016}
M.~L\'opez de Prado,
\emph{Building diversified portfolios that outperform out of sample},
Journal of Portfolio Management \textbf{42}(4) (2016), 59--69.

\bibitem{cotton2022blog}
P.~Cotton,
\emph{Schur complementary portfolios fix hierarchical risk parity},
Geek Culture (Medium), 2022.
\url{https://medium.com/geekculture/schur-complementary-portfolios-fix-hierarchical-risk-parity-28b0efa1f35f}

\bibitem{palomar2025}
D.~P.~Palomar,
\emph{Portfolio Optimization: Theory and Application},
Cambridge University Press, 2025. doi:10.1017/9781009428095.

\bibitem{cotton2024schur}
P.~Cotton,
\emph{Schur complementary allocation: a unification of hierarchical and
optimization-based portfolio construction}, 2024.
arXiv:2411.05807. \url{https://schur.microprediction.org}

\bibitem{antonov2025}
A.~Antonov, A.~Lipton, and M.~L\'opez de Prado,
\emph{Overcoming Markowitz's instability with the help of the hierarchical
risk parity (HRP): theoretical evidence},
in Transactions of ADIA Lab, World Scientific, 2025, pp.~87--121.
doi:10.1142/9789819813049\_0003.

\bibitem{mograby2025}
G.~Mograby,
\emph{Hierarchical minimum variance portfolios: a theoretical and algorithmic
approach}, 2025. arXiv:2503.12328.

\bibitem{ledoitwolf2004}
O.~Ledoit and M.~Wolf,
\emph{A well-conditioned estimator for large-dimensional covariance
matrices},
Journal of Multivariate Analysis \textbf{88}(2) (2004), 365--411.

\end{thebibliography}

\end{document}
